\section{Rezultati}
\label{sec:rezultati}

U ovom poglavlju prikazuju se i interpretiraju rezultati
eksperimentalne evaluacije četiri konfiguracije vektorskih
indeksa opisane u poglavlju~\ref{sec:metodologija}.
Prikazane su deskriptivne statistike latencija, distribucije
vizualizirane dijagramima kutije, memorijski otisak po
konfiguraciji te krivulje kompromisa između tačnosti i
latencije dobivene parametarskim pregledom.
Na osnovu ovih rezultata vrši se formalna verifikacija
istraživačkih hipoteza H1, H2 i H3.

\subsection{Deskriptivna statistika latencija}
\label{sec:stat_latencija}

Za svaku od četiri konfiguracije indeksa izmjerena su
latencija pretrage $1{.}000$ upita s po pet ponavljanja,
što daje ukupno $5{.}000$ mjerenja po konfiguraciji.
Deskriptivni statistički pokazatelji prikazani su u
tabeli~\ref{tab:latencija_stat}, a vizualizacija
distribucija na logaritamskoj skali prikazana je na
slici~\ref{fig:boxplot_latencija}.

\begin{table}[htbp]
\centering
\caption{Deskriptivna statistika latencija pretrage po
         konfiguraciji indeksa. Sve vrijednosti
         izražene su u milisekundama.}
\label{tab:latencija_stat}
\begin{tabular}{lrrrrr}
\hline
\textbf{Konfiguracija} &
\textbf{Mean} &
\textbf{Median} &
\textbf{Std} &
$\mathbf{P_{95}}$ &
$\mathbf{P_{99}}$ \\
\hline
Flat   & 53{,}22 & 51{,}95 &  6{,}54 &  54{,}89 & 100{,}01 \\
IVF    &  8{,}67 &  8{,}65 &  0{,}60 &   9{,}61 &  10{,}01 \\
IVF-PQ &  1{,}12 &  1{,}11 &  0{,}07 &   1{,}24 &   1{,}28 \\
HNSW   &  6{,}95 &  4{,}28 &  9{,}24 &  44{,}08 &  46{,}85 \\
\hline
\end{tabular}
\end{table}

Iz tabele~\ref{tab:latencija_stat} uočava se da
\foreignlanguage{english}{Flat} indeks postiže
aritmetičku sredinu latencije od $53{,}22$\,ms uz
standardnu devijaciju od $6{,}54$\,ms,
što odgovara iscrpnoj linearnoj pretrazi sa složenošću
$\mathcal{O}(n)$ nad $500{.}000$ vektora.
\foreignlanguage{english}{IVF} indeks postiže srednju
latenciju od $8{,}67$\,ms uz standardnu devijaciju od
$0{,}60$\,ms, što ukazuje na visoko konzistentno
ponašanje klasterski ograničene pretrage.
\foreignlanguage{english}{IVF-PQ} je najbrža
konfiguracija s aritmetičkom sredinom od $1{,}12$\,ms
i standardnom devijacijom od $0{,}07$\,ms,
što je posljedica aproksimativnog računanja udaljenosti
nad komprimiranim vektorima direktno u memoriji procesa.

Posebno se ističe \foreignlanguage{english}{HNSW}
konfiguracija, čija je medijana latencije $4{,}28$\,ms,
ali standardna devijacija dostiže $9{,}24$\,ms,
znatno višu vrijednost od svih ostalih konfiguracija.
Percentilne vrijednosti $P_{95} = 44{,}08$\,ms
i $P_{99} = 46{,}85$\,ms otkrivaju prisustvo izrazito
dugog desnog repa distribucije latencija, vidljivog
na slici~\ref{fig:boxplot_latencija}.
Ova pojava pripisuje se arhitekturalnom overheadu
\foreignlanguage{english}{Qdrant} servisa koji radi
kao produkcijski mikroservis s mrežnim interfejsom.
Izmjereni mrežni overhead iznosi $0{,}68$\,ms i
oduzet je na nivou medijane pri apsolutnom poređenju
s \foreignlanguage{english}{FAISS} konfiguracijama,
u skladu s metodologijom opisanom u
poglavlju~\ref{sec:indeksi}.

\begin{figure}[htbp]
    \centering
    \includegraphics[width=0.9\textwidth]{%
        ../images/boxplot_latency.pdf}
    \caption{Raspodjela latencija po konfiguraciji
             indeksa prikazana na logaritamskoj skali.
             Uočava se izrazito dugi desni rep
             distribucije za
             \foreignlanguage{english}{HNSW}
             konfiguraciju u poređenju s kompaktnim
             rasporedima
             \foreignlanguage{english}{IVF} i
             \foreignlanguage{english}{IVF-PQ}.}
    \label{fig:boxplot_latencija}
\end{figure}

\subsection{Memorijski otisak indeksa}
\label{sec:memorija}

Memorijski otisak svake konfiguracije prikazan je na
slici~\ref{fig:memorija}.
\foreignlanguage{english}{Flat} indeks zauzima
$1{.}464{,}8$\,MB jer se svi vektori pohranjuju u
punoj \foreignlanguage{english}{float32} preciznosti.
\foreignlanguage{english}{IVF} zauzima
$1{.}469{,}4$\,MB, blago više od
\foreignlanguage{english}{Flat} indeksa zbog pohrane
centroida klastera i pratećih posting lista.
\foreignlanguage{english}{IVF-PQ} drastično reducira
memorijski otisak na $35{,}8$\,MB zahvaljujući
produktnoj kvantizaciji opisanoj u
poglavlju~\ref{sec:vektorska_kvantizacija},
što iznosi kompresijski omjer od $40{,}88\times$ u
odnosu na \foreignlanguage{english}{Flat} indeks.
\foreignlanguage{english}{HNSW} zauzima
$2{.}139{,}1$\,MB, znatno više od svih ostalih
konfiguracija, jer grafovska struktura pohranjuje
i vektorske vrijednosti i veze grafovskog indeksa
na svakom sloju hijerarhije.

\begin{figure}[htbp]
    \centering
    \includegraphics[width=0.9\textwidth]{%
        ../images/memory_comparison.pdf}
    \caption{Memorijski otisak indeksa po
             konfiguraciji u megabajtima.
             \foreignlanguage{english}{IVF-PQ}
             postiže kompresijski omjer od
             $40{,}88\times$ u odnosu na
             \foreignlanguage{english}{Flat},
             dok \foreignlanguage{english}{HNSW}
             premašuje granicu od $2$\,GB.}
    \label{fig:memorija}
\end{figure}

\subsection{Sweep parametara i krivulje kompromisa}
\label{sec:sweep}

Parametarskim pregledom ispitan je uticaj promjene
ključnih parametara pretrage na kompromis između
tačnosti i latencije za sva tri aproksimativna indeksa.
Za \foreignlanguage{english}{IVF} i
\foreignlanguage{english}{IVF-PQ} variran je parametar
$n_{\mathrm{probe}} \in \{8,\,16,\,32,\,64,\,128\}$,
a za \foreignlanguage{english}{HNSW} parametar
$ef_{\mathrm{search}} \in \{50,\,100,\,200,\,400\}$.
Rezultati su prikazani u
tabeli~\ref{tab:sweep_rezultati} i vizualizirani
na slici~\ref{fig:tradeoff}.

\begin{table}[htbp]
\centering
\caption{Rezultati parametarskog pregleda za tri
         aproksimativna indeksa s variranim
         parametrima pretrage.}
\label{tab:sweep_rezultati}
\begin{tabular}{llrr}
\hline
\textbf{Indeks} &
\textbf{Parametar} &
\textbf{Recall@10} &
\textbf{Latencija (ms)} \\
\hline
IVF &
$n_{\mathrm{probe}} = 8$   & 0{,}8848 &  2{,}25 \\
IVF &
$n_{\mathrm{probe}} = 16$  & 0{,}9307 &  4{,}41 \\
IVF &
$n_{\mathrm{probe}} = 32$  & 0{,}9564 &  8{,}62 \\
IVF &
$n_{\mathrm{probe}} = 64$  & 0{,}9790 & 17{,}11 \\
IVF &
$n_{\mathrm{probe}} = 128$ & 0{,}9945 & 34{,}34 \\
\hline
IVF-PQ &
$n_{\mathrm{probe}} = 8$   & 0{,}6448 & 0{,}31 \\
IVF-PQ &
$n_{\mathrm{probe}} = 16$  & 0{,}6611 & 0{,}57 \\
IVF-PQ &
$n_{\mathrm{probe}} = 32$  & 0{,}6679 & 1{,}08 \\
IVF-PQ &
$n_{\mathrm{probe}} = 64$  & 0{,}6718 & 2{,}11 \\
IVF-PQ &
$n_{\mathrm{probe}} = 128$ & 0{,}6741 & 4{,}19 \\
\hline
HNSW &
$ef_{\mathrm{search}} = 50$  & 0{,}9872 & 6{,}12 \\
HNSW &
$ef_{\mathrm{search}} = 100$ & 0{,}9943 & 6{,}40 \\
HNSW &
$ef_{\mathrm{search}} = 200$ & 0{,}9979 & 7{,}82 \\
HNSW &
$ef_{\mathrm{search}} = 400$ & 0{,}9991 & 9{,}03 \\
\hline
\end{tabular}
\end{table}

Za \foreignlanguage{english}{IVF} indeks uočena je
gotovo linearna veza između parametra
$n_{\mathrm{probe}}$ i latencije pretrage.
Povećanjem $n_{\mathrm{probe}}$ s $8$ na $128$,
latencija raste s $2{,}25$\,ms na $34{,}34$\,ms,
tj.\ oko $15\times$, dok \foreignlanguage{english}%
{recall@10} raste s $0{,}8848$ na $0{,}9945$,
porast od $10{,}97$ procentnih poena.
Time je potvrđeno da $n_{\mathrm{probe}}$ pruža
kontinuiran i predvidiv kompromis između tačnosti
i latencije za \foreignlanguage{english}{IVF}
indeks.

Za \foreignlanguage{english}{IVF-PQ} indeks
primjećuje se fundamentalno drugačije ponašanje.
Povećanjem $n_{\mathrm{probe}}$ s $8$ na $128$,
latencija raste oko $13\times$, ali
\foreignlanguage{english}{recall@10} mijenja se
samo s $0{,}6448$ na $0{,}6741$, porast od svega
$2{,}93$ procentna poena.
Iz ovoga slijedi da je gornja granica tačnosti
\foreignlanguage{english}{IVF-PQ} fundamentalno
određena gubitkom informacija uslijed produktne
kvantizacije, a ne opsegom pretrage klastera.
Šire pretraživanje klastera ne može nadoknaditi
aproksimativnost pri računanju udaljenosti nad
komprimiranim kodovima.

\foreignlanguage{english}{HNSW} indeks pokazuje
najpovoljniji kompromis između tačnosti i latencije
od sva tri razmatrana aproksimativna pristupa.
Već pri $ef_{\mathrm{search}} = 50$ postiže se
\foreignlanguage{english}{recall@10} od $0{,}9872$
uz latenciju od $6{,}12$\,ms. Povećanjem na
$ef_{\mathrm{search}} = 400$ tačnost dostiže
$0{,}9991$ uz latenciju od $9{,}03$\,ms,
što je porast od svega $2{,}91$\,ms uz povećanje
tačnosti za $1{,}19$ procentnih poena.
Ovakvo ponašanje ukazuje na to da grafovska
struktura \foreignlanguage{english}{HNSW} indeksa
efikasno pronalazi gotovo sve relevantne rezultate
čak i s relativno malim skupom kandidata
$ef_{\mathrm{search}}$.

\begin{figure}[htbp]
    \centering
    \includegraphics[width=0.9\textwidth]{%
        ../images/tradeoff_curve.pdf}
    \caption{Krivulje kompromisa između tačnosti i
             latencije za tri aproksimativna indeksa.
             \foreignlanguage{english}{HNSW} i
             \foreignlanguage{english}{IVF} operišu
             u visokom opsegu tačnosti (iznad $0{,}88$),
             dok je \foreignlanguage{english}{IVF-PQ}
             ograničen na
             \foreignlanguage{english}{recall@10}
             između $0{,}64$ i $0{,}68$ bez obzira
             na vrijednost $n_{\mathrm{probe}}$.}
    \label{fig:tradeoff}
\end{figure}

\subsection{Verifikacija hipoteza}
\label{sec:verifikacija_hipoteza}

Na osnovu prikazanih rezultata vrši se formalna
verifikacija istraživačkih hipoteza postavljenih
u poglavlju~\ref{sec:ciljevi}. Zbirni pregled
tačnosti, latencije i memorijskog otiska prikazan
je u tabeli~\ref{tab:zbirni_pregled}.

\begin{table}[htbp]
\centering
\caption{Zbirni pregled ključnih metrika po
         konfiguraciji indeksa.}
\label{tab:zbirni_pregled}
\begin{tabular}{lrrr}
\hline
\textbf{Konfiguracija} &
\textbf{Recall@10} &
\textbf{Mean (ms)} &
\textbf{Memorija (MB)} \\
\hline
Flat   & 1{,}0000 & 53{,}22 & 1{.}464{,}8 \\
IVF    & 0{,}9564 &  8{,}67 & 1{.}469{,}4 \\
IVF-PQ & 0{,}6679 &  1{,}12 &      35{,}8 \\
HNSW   & 0{,}9961 &  6{,}95 & 2{.}139{,}1 \\
\hline
\end{tabular}
\end{table}

\subsubsection{Hipoteza H1}
\label{sec:h1}

Hipoteza H1 tvrdi da \foreignlanguage{english}{HNSW}
indeks postiže statistički značajno nižu latenciju
upita u poređenju s \foreignlanguage{english}{IVF}
indeksom uz zadržavanje
\foreignlanguage{english}{recall@10} iznad $90\%$.

Wilcoxon \foreignlanguage{english}{signed-rank} test
primijenjen je na parove izmjerenih latencija
\foreignlanguage{english}{HNSW} i
\foreignlanguage{english}{IVF} konfiguracija.
Testna statistika iznosi $W = 51{.}069$ uz
vrijednost $p \approx 0$, što je daleko ispod praga
$p < 0{,}05$. Nulta hipoteza o jednakim
distribucijama latencija odbacuje se, a razlika
između konfiguracija proglašava se statistički
značajnom.
Medijana latencije \foreignlanguage{english}{HNSW}
iznosi $4{,}28$\,ms naspram $8{,}65$\,ms za
\foreignlanguage{english}{IVF}, a izmjereni
\foreignlanguage{english}{recall@10} od $0{,}9961$
premašuje prag od $90\%$.
Oba uvjeta hipoteze H1 su ispunjena.

\textbf{Hipoteza H1 je potvrđena.}

Ipak, napominje se važna praktična napomena.
Medijana ne opisuje potpuno ponašanje
\foreignlanguage{english}{HNSW} indeksa u uvjetima
opterećenja. Vrijednosti $P_{95} = 44{,}08$\,ms i
$P_{99} = 46{,}85$\,ms ukazuju na to da se u
najlošijih $5\%$ slučajeva latencija
\foreignlanguage{english}{HNSW} povećava više od
deset puta u odnosu na medijanu, što je znatno
nepovoljnije od \foreignlanguage{english}{IVF}
($P_{95} = 9{,}61$\,ms).
Ovaj nalaz pripisuje se mrežnom i
serijalizacijskom overheadu
\foreignlanguage{english}{Qdrant} mikroservisa,
koji u određenim scenarijima uzrokuje povećanu
varijabilnost odgovora. U produkcijskim sistemima
osjetljivim na latenciju repova, ova razlika
je od kritičnog inžinjerskog značaja i treba
je uzeti u obzir pri odabiru indeksne strukture.

\subsubsection{Hipoteza H2}
\label{sec:h2}

Hipoteza H2 tvrdi da primjena produktne
kvantizacije smanjuje memorijski otisak za
najmanje $4\times$ u odnosu na nekomprimirani
\foreignlanguage{english}{flat} indeks, uz
smanjenje \foreignlanguage{english}{recall@10}
manje od $5$ procentnih poena.

Izmjereni memorijski otisak
\foreignlanguage{english}{Flat} indeksa iznosi
$1{.}464{,}8$\,MB, a
\foreignlanguage{english}{IVF-PQ} indeksa
$35{,}8$\,MB, što daje kompresijski omjer od
$40{,}88\times$. Ovaj omjer znatno premašuje
prag od $4\times$ i prvi uvjet hipoteze H2
ispunjen je s rezervom.

Međutim, izmjereni \foreignlanguage{english}%
{recall@10} za \foreignlanguage{english}{IVF-PQ}
iznosi $0{,}6679$, dok gornju granicu tačnosti
definira \foreignlanguage{english}{Flat} indeks
s vrijednošću $1{,}0000$.
Pad tačnosti iznosi $0{,}3321$, odnosno
$33{,}21$ procentnih poena, što daleko premašuje
propisani prag od $5$ procentnih poena.
Drugi uvjet hipoteze H2 nije ispunjen.

\textbf{Hipoteza H2 nije potvrđena u cjelosti.}

Memorijska kompresija od $40{,}88\times$
potvrđena je s značajnom rezervom, ali pad
tačnosti od $33{,}21$ procentnih poena ukazuje
na fundamentalno ograničenje produktne
kvantizacije u opisanoj konfiguraciji.
Ovo je direktna posljedica agresivnih parametara
kompresije $M = 64$ podvektora i
$n_{\mathrm{bits}} = 8$ bita po podvektoru.
Pri ovakvoj konfiguraciji vektori su podijeljeni
na $64$ disjunktna podprostora od kojih svaki
bilježi grešku aproksimacije centroidu
odgovarajućeg \foreignlanguage{english}{codebook}
mehanizma. Kvantizacijska greška u svakom od
$64$ potprostora akumulira se pri rekonstrukciji
aproksimativnih udaljenosti prema
jednadžbi~\ref{eq:pq_dist}, što degradira
ukupnu tačnost pretrage.
Iz parametarskog pregleda vidljivo je da
povećanje $n_{\mathrm{probe}}$ s $8$ na $128$
donosi svega $2{,}93$ procentna poena poboljšanja
tačnosti uz $13\times$ duži odziv,
što potvrđuje da je gornja granica tačnosti
određena kvantizacijom, a ne opsegom pretrage.
Smanjenjem broja potprostora $M$ ili povećanjem
$n_{\mathrm{bits}}$ mogao bi se postići povoljniji
balans između kompresije i tačnosti na štetu
memorijske uštedine.

\subsubsection{Hipoteza H3}
\label{sec:h3}

Hipoteza H3 tvrdi da kombinacija
\foreignlanguage{english}{IVF} i produktne
kvantizacije pruža bolji kompromis između
tačnosti i latencije od samostalnog
\foreignlanguage{english}{HNSW} pristupa
kada je memorijski budžet ograničen ispod $2$\,GB.

Izmjereni memorijski otisak
\foreignlanguage{english}{IVF-PQ} indeksa iznosi
$35{,}8$\,MB, što je znatno ispod postavljenog
budžeta od $2$\,GB. Nasuprot tome, memorijski
otisak \foreignlanguage{english}{HNSW} indeksa
iznosi $2{.}139{,}1$\,MB i premašuje zadani
budžet za $91{,}1$\,MB. U uvjetima strogog
memorijskog ograničenja od $2$\,GB,
\foreignlanguage{english}{HNSW} ne može biti
primijenjen, dok je \foreignlanguage{english}%
{IVF-PQ} jedina razmatrana konfiguracija koja
staje u zadani budžet.

\textbf{Hipoteza H3 je potvrđena.}

\foreignlanguage{english}{IVF-PQ} predstavlja
jedini izvedivi kompromis između tačnosti i
latencije unutar budžeta od $2$\,GB.
Iz krivulje kompromisa na
slici~\ref{fig:tradeoff} vidljivo je da
\foreignlanguage{english}{IVF-PQ} nudi
latencije između $0{,}31$\,ms i $4{,}19$\,ms
uz tačnost u rasponu $0{,}64$ do $0{,}67$,
dok je \foreignlanguage{english}{HNSW}
isključen iz razmatranja zbog nedostupnosti
unutar budžeta. Napominje se da ovaj zaključak
važi isključivo pod postavljenim memorijskim
ograničenjem. U scenarijima bez strogog
budžeta, \foreignlanguage{english}{HNSW}
ostaje superiorniji izbor s obzirom na
znatno višu tačnost ($0{,}9961$ naspram
$0{,}6679$) uz usporedivu ili nižu latenciju
na nivou medijane.
